\section{The Proposed Framework} \label{sec:33-research01-framework}

\subsection{Problem statement}
\begin{foldable}
    Given a \textit{black-box} pre-trained teacher network $T$ and a small set of labeled images $\mathcal{D}=\{x_{i},y_{i}\}_{i=1}^{N}$, our goal is to train a student network $S$ on $\mathcal{D}$ such that $S$ can approximate the performance of $T$.
\end{foldable}

\begin{foldable}
    A direct solution is to use standard KD \cite{hinton2015distilling}, which trains $S$ to match both the ground-truth labels $y_{i}$ and the teacher's predictions $T(x_{i})$ via minimizing:
\end{foldable}
\begin{equation}
    \mathcal{L}_{\text{KD}}=\mathbb{E}_{(x_{i},y_{i})\sim{\cal D}}\Bigl[(1-\lambda)\mathcal{L}_{\text{CE}}(S(x_{i}),y_{i})+\lambda\mathcal{L}_{\text{KL}}(S(x_{i}),T(x_{i}))\Bigr]\label{ch3:eq:StandardKD}
\end{equation}
\begin{foldable}
    where $\mathcal{L}_{\text{CE}}$ and $\mathcal{L}_{\text{KL}}$ are the cross-entropy and the Kullback-Leibler divergence losses, $S(x_{i})$ and $T(x_{i})$ are the student\textquoteright s and teacher's predictions (i.e. \textit{class probabilities}), and $\lambda$ is a coefficient to balance the losses.
    Note that the \textit{temperature} factor in the original paper is not used as it requires access to the teacher's logits, which violates our ``black-box teacher'' assumption.
    Moreover, with a small distillation set ${\cal D}$, common KD approaches lead to a \textit{sub-optimal} model as they require lots of training images \cite{hinton2015distilling,kim2018paraphrasing,ahn2019variational,tian2020contrastive}.
\end{foldable}

\subsection{Proposed method\label{ch3:subsec:Proposed-method}}
\begin{foldable}
    We propose our novel method \textbf{DivBFKD} to perform \textit{black-box few-shot} KD.
    Our method has two phases: \textit{generation} and \textit{distillation}.
    In the generation phase, we train a WGAN to produce synthetic images.
    In the distillation phase, we train the student with the combined real and synthetic images.
\end{foldable}

\subsubsection{Generation phase.}
\begin{foldable}
    We train a WGAN model \cite{gulrajani2017improved} with several key changes to generate diverse synthetic images.
    Like standard WGAN, our WGAN consists of a generator $G$ and a discriminator $D$.
    Iteratively, $G$ produces a synthetic image \textbf{$\tilde{x}=G\left(z\right)$} from a random latent vector $\ensuremath{z\sim\mathcal{N}\left(0,I\right)}$.
    Meanwhile, $D$ predicts a scalar score indicating the likelihood of an input image being real.
    Jointly trained with opposite objectives, $G$ tries to fool $D$ by producing realistic images while $D$ tries to best distinguish between real images $x$ and synthetic images $\tilde{x}$.
    These objectives can be formulated into the generator loss ${\cal L}_{G}$ and the discriminator loss ${\cal L}_{D}$:
\end{foldable}
\begin{equation}
    \mathcal{L}_{G}=-\mathbb{E}_{z\sim{\cal N}(0,I)}[D(G(z))],\label{ch3:eq:Loss-G-WGAN}
\end{equation}
\begin{equation}
    \mathcal{L}_{D}=\mathbb{E}_{z\sim{\cal N}(0,I)}[D(G(z))]-\mathbb{E}_{x\sim{\cal D}}[D(x)].\label{ch3:eq:Loss-D-WGAN}
\end{equation}
\begin{foldable}
    However, if we only use the few-shot set $\mathcal{D}$ to train the WGAN, $G$ will generate synthetic images similar to the images in $\mathcal{D}$ \cite{karras2020training,yang2021data}, which are not diverse enough to facilitate effective distillation.
    To improve diversity of images generated by $G$, we propose the use of synthetic images that the teacher $T$ can predict confidently.
    We refer to these images as \textit{high-confidence images}.
\end{foldable}

\begin{foldable}
    \textbf{High-confidence images.} Given a synthetically generated (fake) image $\tilde{x}$, we define its \textbf{\textit{confidence-score}} as $c_{\tilde{x}}=\max{k\in\{0,...,K-1\}}T(\tilde{x})$, by taking the maximum over $T$'s predictive probability vector.
    For example, given three classes $\{0,1,2\}$ and class probabilities $T(\tilde{x})=[0.1,0.7,0.2]$, then $c_{\tilde{x}}=0.7$.
    We define $\tilde{x}$ a \textbf{\textit{high-confidence image}} if $c_{\tilde{x}}$ is above a \textit{confidence-threshold} $\tau\in[0,1]$ i.e. $c_{\tilde{x}}\geq\tau$.
    Intuitively, when $c_{\tilde{x}}$ is high, $\tilde{x}$ is likely close to the teacher's training images as it can be confidently classified.
\end{foldable}

\begin{foldable}
    \textbf{Adaptive thresholds.} The critical factor in determining a high-confidence image is the confidence-threshold $\tau$.
    We observe that the teacher has a \textit{class-specific bias} i.e. it predicts different classes with different confidence levels.
    Take the dataset Fashion (Fig. \ref{ch3:fig:Pre-clamp-tau}) for example, the averaged confidence-score is $>0.99$ for classes ``trouser'' and ``bag'', but is only $0.75$ for class ``shirt''.
    In this case, fixing a single value $\tau=0.95$ may lead to \textit{high-confidence images biased to certain classes} such as ``trouser'' and ``bags'' but may be \textit{extremely demanding} for the class ``shirt''.
    Thus, a single $\tau$ for all classes will make high-confidence images imbalanced.
    Further, considering any single class, the confidence-scores have a long-tailed distribution, where they concentrate near 1 and decay with lower values.
    Even though $\tau$ can be simply set to the averaged confidence-score, it would provide little flexibility and interpretability and also make \textit{high-confidence images biased towards outliers}.
\end{foldable}

\begin{figure}
    \centering
    \includegraphics[viewport=0bp 20bp 455bp 174bp,clip,width=0.75\columnwidth]{./figures/30-research01/33-research01-framework-preclamptau.pdf}
    \caption{Distributions of confidence-scores across classes on Fashion. They concentrate around 1 and have a long-tail. Red bars show adaptive confidence-thresholds $\tau$.}
    \label{ch3:fig:Pre-clamp-tau}
\end{figure}

\begin{foldable}
    To address these problems, we use \textit{adaptive thresholds} $\{\tau^{k}\}_{k=0}^{K-1}$ that is class-specific based on the $q$-quantile of the confidence-scores of real images.
    Specifically, for each of the $K$ classes:
\end{foldable}
\begin{equation}
    \tau^{k}=\text{quantile}_{q}(\{c_{x_{i}}\mid(x_{i},y_{i})\in\mathcal{D}\text{\text{ and }}y_{i}=k\}_{i=1}^{N}),q\in\left[0,1\right].\label{ch3:eq:Adaptive-threshold}
\end{equation}
\begin{foldable}
    We choose $q=0.1$ i.e. high-confidence images must have teacher's confidence higher than at least 10\% of those for the real ones, as perceived by the teacher $T$.
    This way, $\tau$ adaptively adjusts to any dataset and any class-wise bias from the teacher.
    Fig. \ref{ch3:fig:Pre-clamp-tau} illustrates $\tau$ as red bars using $q=0.1$ on Fashion.
\end{foldable}

\begin{foldable}
    \textbf{Our WGAN training scheme.} We aim to harness the teacher's guidance to promote diversity of synthetic images.
    Fundamentally, synthetic images would be more diverse if the WGAN can learn on images that resemble the \textbf{teacher's unknown training images} $\mathcal{D^{*}}$, whose diversity is superior to few-shot images $\mathcal{D}$.
    While there is \textit{no direct access} to $\mathcal{D^{*}}$, high-confidence images can serve as a satisfactory alternative in terms of both quantity (new high-confidence images are generated at every WGAN training step) and quality (they likely reside close to $\mathcal{D}^{*}$, as we previously discussed).
    In other words, high-confidence images implicitly allow our WGAN to learn from $\mathcal{D}^{*}$.
\end{foldable}

\begin{figure}
    \centering
    \includegraphics[width=1\columnwidth]{./figures/30-research01/33-research01-framework-teaser.pdf}
    \caption{Illustration of our method DivBFKD. \textbf{(a)} \textbf{Generation:} We train our WGAN with the losses ${\cal L}_{G}$ and ${\cal L}_{D}^{\text{new}}$. When optimizing $G$, we construct the \textit{high-confidence set} ${\cal H}_{n}=\left\{\tilde{x}=G(z)\mid c_{\tilde{x}}\protect\geq\tau^{\tilde{y}}\right\} $ from synthetic images $\tilde{x}$, with their confidence-score $c_{\tilde{x}}$, pseudo-label $\tilde{y}$, and adaptive threshold $\tau^{\tilde{y}}$. When optimizing $D$, we sample $x$ from the combined real and high-confidence images set $\mathcal{D}\cup{\cal H}_{n-1}$. \textbf{(b) Distillation:} We augment few images in $\mathcal{D}$ with synthetic images in ${\cal D}_{G}$ to train the student $S$ on ${\cal D}\cup{\cal D}_{G}$.}
    \label{ch3:fig:Teaser}
\end{figure}

\begin{foldable}
    By this motivation, we introduce high-confidence images in the role of real images to the adversarial learning (Fig. \ref{ch3:fig:Teaser}(a)), exposing our WGAN to a wider variety of training images.
    At each $n$-th training step:
\end{foldable}
\begin{itemize}
    \item The generator $G$ is trained with Eq. \ref{ch3:eq:Loss-G-WGAN}. For a fake image $\tilde{x}$, we compute its corresponding confidence-score $c_{\tilde{x}}$ and pseudo-label $\tilde{y}=\argmax{k\in\{0,...,K-1\}}T(\tilde{x})$. Using the adaptive thresholds $\{\tau^{k}\}_{k=0}^{K-1}$ computed with Eq. \ref{ch3:eq:Adaptive-threshold}, we add qualified $\tilde{x}$ to the \textit{high-confidence set} ${\cal H}_{n}=\{\tilde{x}\mid c_{\tilde{x}}\geq\tau^{\tilde{y}}\}$. Then, ${\cal H}_{n}$ is stored to train the discriminator $D$ in the next step.
    \item Training the discriminator $D$ requires both real and fake images. However, our `real' images include real images that actually come from ${\cal D}$ and high-confidence images that come from the \textit{high-confidence set} ${\cal H}_{n-1}$. Note that ${\cal H}_{n-1}$ is constructed in the previous step, and ${\cal H}_{0}=\emptyset$. With these changes, we train our discriminator with the new loss as:
\end{itemize}

\begin{equation}
    \mathcal{L}_{D}^{\textnormal{new}}=\mathbb{E}_{z\sim{\cal N}(0,I)}[D(G(z))]-\mathbb{E}_{x\sim\left({\cal D}\cup{\cal H}_{n-1}\right)}[D(x)].\label{ch3:eq:Loss-D-DivBFKD}
\end{equation}

\begin{foldable}
    Viewing as an adversarial game, when $D$ learns to appreciate high-confidence synthetic images as realistic, it also encourages $G$ to create more synthetic images like those.
    This interaction forms a positive feedback loop that helps $G$ to be more diverse.
    Furthermore, we emphasize that forming ${\cal H}_{n}$ is a \textit{cheap} operator and requires small memory.
    \textit{These modifications not only effectively improve the diversity of synthetic images, but also incur only trivial computation overhead.}
\end{foldable}

\begin{foldable}
    \textbf{Summary of our novelty in WGAN.} We introduce a \textit{new training scheme} for WGAN using \textit{high-confidence images} based on the teacher-supervised \textit{adaptive thresholds}.
    Since we aim to generate images close to the teacher's training images, our synthetic images are more \textit{diverse} and better cover \textit{unseen images}, significantly improving the student's generalization, as shown in Section \ref{sec:34-research01-experiments}.
\end{foldable}

\begin{algorithm}[ht]
    \SetKwInOut{KwIn}{Input}
    \SetKwInOut{KwOut}{Output}
    \SetKwComment{Comment}{}{}

    \caption{Proposed method \textbf{DivBFKD}}
    \label{ch3:alg:Proposed-method-DivBFKD}

    \KwIn{teacher $T$; labeled set $\mathcal{D}=\{x_{i},y_{i}\}_{i=1}^{N}$; quantile $q$; synthetic budget $M$; batch size $B$}
    \KwOut{student network $S$}

    \vspace{-0.5\baselineskip}
    \Comment{\noindent\hrulefill}

    \vspace{0.5\baselineskip}
    Initialize generator $G$, discriminator $D$, student $S$, and high-confidence set $\mathcal{H}_0 \leftarrow \emptyset$ \\
    \Comment{$\triangleright$ \textbf{\textit{Stage 1: Generation}}}
    Infer $T$ on $\mathcal{D}$ to compute adaptive thresholds $\{\tau^k\}_{k=1}^K$ with $q$ \hfill via Eq.~\eqref{ch3:eq:Adaptive-threshold} \\
    \While{$G$ not converged}{
        Sample latent vectors $\{z_i\}_{i=1}^B \sim \mathcal{N}(0,I)$ \\
        Generate synthetic images $\{\tilde{x}_i\}_{i=1}^B \leftarrow G(z_i)$ \\
        Sample $\{x_i\}_{i=1}^B \sim (\mathcal{D} \cup \mathcal{H}_n)$ as real images \\
        Update $G$ on $\left\{\tilde{x}_i\right\}_{i=1}^B$ via $\mathcal{L}_{G}$ in Eq. \ref{ch3:eq:Loss-G-WGAN} \\
        Update $D$ on $\left\{x_i\right\}_{i=1}^B$ and $\left\{\tilde{x}_i\right\}_{i=1}^B$ via $\mathcal{L}_{D}^{new}$ in Eq. \ref{ch3:eq:Loss-D-DivBFKD} \\
        Compute pseudo-labels $\{\tilde{y}_i\}$ and confidence scores $\{c_{\tilde{x}_i}\}$ for $\{\tilde{x}_i\}$ \\
        Construct $\mathcal{H}_n=\left\{\tilde{x_{i}} \mid c_{\tilde{x}_i}\ge\tau^{\tilde{y}_i}\right\}_{i=1}^B$ \\
    }

    \vspace{0.5\baselineskip}
    \Comment{$\triangleright$ \textbf{\textit{Stage 2: Distillation}}}
    Construct $\mathcal{D}_{\text{KD}}$ with $N$ real and $M$ synthetic images \hfill via Eq.~\eqref{ch3:eq:Distillation-set} \\
    \While{$S$ not converged}{
        Update $S$ on $\mathcal{D}_{\text{KD}}$ via $\mathcal{L}_{\text{KD}}$ \hfill via Eq.~\eqref{ch3:eq:Loss-S}
    }

    \vspace{0.5\baselineskip}
    \Return{$S$}
\end{algorithm}

\subsubsection{Distillation phase.}
\begin{foldable}
    After training WGAN in the generation phase, we use it to generate synthetic images.
    We employ our generator $G$ to build a set of synthetic images ${\cal D}_{G}$, and their pseudo-labels (i.e. class probabilities) are obtained via the teacher $T$.
    Following BBKD \cite{wang2020neural} and FS-BBT \cite{nguyen2022black}, we construct a \textit{distillation set} $\mathcal{D}_{\text{KD}}$ comprising of $N$ real images from $\mathcal{D}$ and $M$ synthetic images from ${\cal D}_{G}$:
\end{foldable}
\begin{equation}
    \mathcal{D}_{\text{KD}}=\left\{x_{i}\in\mathcal{D}\right\} _{i=1}^{N}\cup\left\{G\left(z_{j}\right)|z_{j}\sim\mathcal{N}\left(0,I\right)\right\} _{j=1}^{M},
    \label{ch3:eq:Distillation-set}
\end{equation}
\begin{foldable}
    where $x_{i}\in{\cal D}$ and $\tilde{x}_{j}\in{\cal D}_{G}$ are real and synthetic images.
    Like other few-shot KD methods \cite{wang2020neural,nguyen2022black}, we train the student $S$ with ${\cal D}_{\text{KD}}$ using a cross-entropy loss:
\end{foldable}
\begin{equation}
    {\cal {\cal L}}_{\text{KD}}=\mathbb{E}_{x_{i}\sim{\cal D}_{\text{KD}}}{\cal L}_{\text{CE}}\left(S\left(x_{i}\right),T\left(x_{i}\right)\right).
    \label{ch3:eq:Loss-S}
\end{equation}

\begin{foldable}
    \textbf{Class balancing.} Inspired by \cite{chen2019data,zhang2022ideal}, we generate a balanced set of synthetic images.
    Recalling that $M$ is the budget of synthetic images for KD and $K$ is the number of classes, each class would have $\frac{M}{K}$ synthetic images.
    We ensure class balance through \textit{rejection sampling} until we achieve the sufficient amount.
\end{foldable}

\begin{foldable}
    We summarize our method in Algorithm \ref{ch3:alg:Proposed-method-DivBFKD}.
\end{foldable}
